% ============================================================================
\clearpage
\section*{v4.10.0-r46 Source-Only: Frame, Axis and Orientation as Readout Data (Branch B)}
\addcontentsline{toc}{section}{v4.10.0-r46 Source-Only: The Readout Branch of the Mediator}
\label{sec:v4100-r46-readout-branch}
% ============================================================================

\paragraph{Purpose.}
Layer v4.10.0-r45 found that the recurring candidates for the OP5 mediator are four ways of supplying a frame \(u\), an axis \(n\) and a time
orientation, and that r35--r44 placed these on the field side (branch H) without saying so.  At the author's choice this layer tests the readout
side (branch B): \(u\) and \(n\) are the rest frame and measuring axis of the detector.  The strict criterion of r45 applies: the branch counts only if the
number of set inputs decreases without a new one and if it predicts something beyond \(N_m\).  It is an exploration with a method fixed before the
computation; it is partly not blind, since the quadratic appearance of \(u\) and the independence of the leading static term from \(r_m\) were
foreseeable by hand.  Earlier layers are not rewritten.

\noindent\textbf{Source-only status.}
This v4.10.0-r46 layer is source-only.  It does not build a PDF, does not assign a new DOI,
and does not authorize a public release.  The v5.3.0 Version DOI
\href{https://doi.org/10.5281/zenodo.24434403}{10.5281/zenodo.24434403} and the Concept DOI
\href{https://doi.org/10.5281/zenodo.19210728}{10.5281/zenodo.19210728} are unchanged by it.

\subsection*{v4.10.0-r46 setting}
\label{subsec:v4100-r46-setting}

The coupling is written covariantly as \(\mathcal L=-h\,\Pi^{\mu\nu}u^\alpha u^\beta R^{(1)}_{\mu\alpha\nu\beta}[\gamma]\) with \(\Pi^{\mu\nu}=-\eta^{\mu\nu}+u^\mu u^\nu-n^\mu n^\nu\), the projector
transverse to \(n\) in the rest frame of \(u\), and the linearized curvature in the convention of the bifocal notes.  The source \(S^{\mu\nu}\) is the coefficient
of \(\tfrac12\gamma_{\mu\nu}h\), and the mediator numerator is the Lorentz-covariant Fierz--Pauli contraction \(S_{\mu\nu}S^{\mu\nu}-\tfrac13(S^\mu{}_\mu)^2\).

\subsection*{v4.10.0-r46 results}
\label{subsec:v4100-r46-results}

\paragraph{Theoremlet v4.10.0-r46.1 (the directed coupling is a transverse tidal readout).}
In the rest frame \(u=(1,0,0,0)\), \(n=e_z\), the source of \(\mathcal L\) equals the hand-built source of the bifocal notes exactly, and its numerator is \(N_m\).
Identically \(\Pi^{\mu\nu}u^\alpha u^\beta R_{\mu\alpha\nu\beta}=-R_{uu}-E_{nn}\) with \(E_{nn}=n^\mu u^\alpha n^\nu u^\beta R_{\mu\alpha\nu\beta}\); both partial sources are conserved, with numerators
\(\tfrac23(a_\perp+k_z^2)^2\) and \(\tfrac23(k_z^2-x)^2\) and a cross term.  \(N_m\) arises from the whole transverse readout, not from one part.
\emph{Proof.} Exact symbolic evaluation (script \texttt{readout\_\allowbreak branch.wl}).  \(\square\)

\paragraph{Theoremlet v4.10.0-r46.2 (orientation and slowed propagation drop out).}
The source is invariant under \(u\mapsto-u\) and the numerator under \(\omega\mapsto-\omega\), so the time orientation of the readout does not enter.  The static
term \(g_m^2\cdot\tfrac23a_\perp^2/(m^2+r_mq)=\tfrac23g_m^2a_\perp^2/m^2-\tfrac23g_m^2a_\perp^2r_mq/m^4+\dots\) has a leading part independent of \(r_m\); with \(r_m=1\) and
\(g_m^2=\tfrac3{100}\) both conditions of r38.2 hold.
\emph{Proof.} As r46.1; series in \(q\).  \(\square\)

\paragraph{Theoremlet v4.10.0-r46.3 (dependence on the motion of the detector).}
For a static source and a detector with velocity \(\beta\), the numerator is \(\tfrac23a_\perp(a_\perp-\beta^2k_z^2)+O(\beta^4)\) for motion along the axis and
\(\tfrac23a_\perp^2+\tfrac23a_\perp^2\beta^2\cos^2\varphi+O(\beta^4)\) for motion across it, \(\varphi\) the azimuth of \(k_\perp\) relative to the motion.
\emph{Proof.} Boost of \(u\) and \(n\) with \(u\cdot u=1\), \(n\cdot n=-1\), \(u\cdot n=0\); exact evaluation and series.  \(\square\)

\paragraph{Remark v4.10.0-r46.4 (reading).}
Against the strict criterion fixed beforehand the branch \emph{carries}.  The set inputs decrease from the directed coupling with frame and axis, the
slowed propagation \(r_m\), the values \(r_m,g_m,m\) and the open time orientation to the existence of a massive spin-2 mediator of mass \(m\) and its coupling
\(g_m\) to the Higgs radius through the readout; \(u\) and \(n\) are data of the measurement, and r32.4 already calls the contrast a tidal readout of a ray along
the axis.  The distinguished frame required in r36 is the detector; no medium, string or seed flow is needed for it, and the questions of r39--r44
do not arise in this branch.  The branch predicts a dependence of the contrast on the motion of the detector relative to the preparation, of relative
size \(\beta^2\) and with the angular structure of r46.3, where branch H would tie it to the motion of a medium; for laboratory velocities the effect is
very small.  Since a massive field with \(r_m=1\) has a group velocity below \(c\), the slowed propagation \(r_m<1\) is not a separate requirement for the leading
contrast (a reading of r33.3).  The remaining field-side question is now stated without frame or orientation: whether the QSC model contains a
massive spin-2 degree of freedom coupling linearly to the Higgs radius; W1 and W2 do not (r33, r34).

\subsection*{v4.10.0-r46 frames}
\label{subsec:v4100-r46-frames}

All quantities are in the matter frame M; \(u\) is the rest frame of the detector and \(n\) its measuring axis.  Time reversal acts as \(u\mapsto-u\) or \(\omega\mapsto-\omega\); the
tachyon map is not used.  No quantity of another frame is compared.

\subsection*{v4.10.0-r46 register entries}
\label{subsec:v4100-r46-register}

\begin{center}
\small
\begin{tabular}{@{}>{\raggedright\arraybackslash}p{0.34\textwidth}>{\raggedright\arraybackslash}p{0.58\textwidth}@{}}
\toprule
statement & status after v4.10.0-r46\\
\midrule
directed coupling & transverse tidal readout of a detector (\(u\), \(n\)) (r46.1)\\
time orientation, slowed propagation & not needed for the leading contrast (r46.2)\\
prediction & dependence on detector motion of order \(\beta^2\) (r46.3)\\
set inputs & massive spin-2 mediator (\(m\)) and its coupling \(g_m\)\\
OP5 & open\\
\bottomrule
\end{tabular}
\end{center}

\subsection*{v4.10.0-r46 notes}
\label{subsec:v4100-r46-notes}

\texttt{notes/\allowbreak research/\allowbreak op5\_\allowbreak readout\_\allowbreak branch\_\allowbreak 2026\_\allowbreak 11\_\allowbreak 23/} (method note commit bb7274a1,
script \texttt{readout\_\allowbreak branch.wl}, output in \texttt{run/}).

\subsection*{v4.10.0-r46 ledger and guard update}
\label{subsec:v4100-r46-ledger-guard}

\[
R_{\rm chk}^{v4100\text{-}r46}=(1,1,0)
\]
for (status recorded; computations checked in the cited script; the existence of the massive spin-2 mediator and its coupling not derived).

The layer spends no guard:
\[
R_{\rm glob}^{(v4100\text{-}r46,+)}=(0,1,1,1,1),
\qquad
\Gamma_{\rm guard}^{v4100\text{-}r46}=(1,1,1,1,1),
\qquad
\beta_4=1/256\ \text{unspent}.
\]
The layer does not imply exact finite \(C5\) zero closure, exact all-readout terminality, \(C6/T2\), a global photon-clock/RH theorem,
or global all-readout closure.  OP1--OP5, C5 and C6/T2 remain open.
